\documentclass[12pt]{article}

% --- Языки и шрифты для XeLaTeX/LuaLaTeX ---
\usepackage{fontspec}
\usepackage{polyglossia}
\setdefaultlanguage{russian}
\setotherlanguage{english}

% Добавляет смайлики :)
\usepackage{wasysym}

% Подходящие шрифты с кириллицей (CMU есть почти везде)
\setmainfont{CMU Serif}
\setsansfont{CMU Sans Serif}
\setmonofont{CMU Typewriter Text}

% Математика и вёрстка
\usepackage{amsmath,amssymb}
\usepackage{geometry}
\geometry{a4paper,margin=2.5cm}

% Ссылки
\usepackage{hyperref}
\hypersetup{
  colorlinks=true,
  linkcolor=blue,
  urlcolor=blue,
  citecolor=blue
}

\begin{document}

\begin{center}
  {\LARGE Домашнее задание \#2}\\[4pt]
  {\large Курс: Введение в маневрирование КА}\\[2pt]
  {\small Дедлайн: \underline{20.10.2025, 23:59} }
\end{center}

\bigskip

\section*{Инструкция}
Для выполнения задания рекомендуется скачать на свой компьютер \href{https://github.com/Chelovechek-move/maneuvering}{этот репозиторий}. В файле \texttt{README} вы найдёте короткие инструкции для начала работы с проектом. В этот раз мы будем работать с одним Jupyter Notebook файлом. Его можно найти тут:
\begin{center}
    \texttt{examples/one\_impulse\_maneuvers\_lecture.ipynb}    
\end{center}
В нём пошагово решены все задачки, которые мы разбирали на лекции. Рекомендуется ещё раз пролистать презентацию и посмотреть, как были получены все результаты.

\section*{Задания}

\begin{enumerate}
    \item \textbf{Одноимпульсный компланарный переход.} \\
    \underline{Дано:} \\
    \underline{Начальная орбита:}
      \begin{equation*}
          \begin{aligned}
            &a_1 = 11\;000~\text{км},\\
            &e_1 = 0.10,\\
            &\omega_1 = 23^{\circ}.
          \end{aligned}
      \end{equation*}
      \medskip
      \underline{Целевая орбита:}
      \begin{equation*}
          \begin{aligned}
            &a_2 = 11\;500~\text{км},\\
            &e_2 = 0.75,\\
            &\omega_2 = 87^{\circ}.
          \end{aligned}
      \end{equation*}
      \underline{Определить:}
       \begin{enumerate}
        \item Точки пересечения орбит;
        \item Импульс скорости $\Delta V$, необходимый для осуществления перехода в соответствующей точке.
      \end{enumerate}

    \item \textbf{Одноимпульсный компланарный переход.} (ещё один \smiley) \\
    \underline{Дано:} \\
    \underline{Начальная орбита:}
      \begin{equation*}
          \begin{aligned}
            &a_1 = 11\;000~\text{км},\\
            &e_1 = 0.10,\\
            &\omega_1 = 23^{\circ}.
          \end{aligned}
      \end{equation*}
      \medskip
      \underline{Целевая орбита:}
      \begin{equation*}
          \begin{aligned}
            &a_2 = 15\;500~\text{км},\\
            &e_2 = 0.10,\\
            &\omega_2 = 87^{\circ}.
          \end{aligned}
      \end{equation*}
      \underline{Определить:}
       \begin{enumerate}
        \item Точки пересечения орбит;
        \item Импульс скорости $\Delta V$, необходимый для осуществления перехода в соответствующей точке.
      \end{enumerate}

      \item \textbf{Коррекция $e$ и $\omega$.} \\
      \underline{Дано:} \\
      \underline{Начальная орбита:}
      \begin{equation*}
          \begin{array}{r@{\;}c@{\;}l @{\qquad} r@{\;}c@{\;}l}
            a_1 & = & 26\;000~\text{км} & \Omega_1 & = & 0^\circ \\
            e_1 & = & 0.7               & \omega_1 & = & 250^\circ \\
            i_1 & = & 64^\circ          & \nu_1    & = & 110^\circ \\
          \end{array}
      \end{equation*}
      \underline{Целевые элементы}
      \begin{align*}
            &e_2 = 0.6 \\
            &\omega_2 = 270^\circ
      \end{align*}
      \underline{Определить:}
       \begin{enumerate}
        \item Импульс скорости $\Delta V$, необходимый для осуществления коррекции.
      \end{enumerate}
      \underline{Построить} график зависимости хар. скорости манёвра коррекции $e$ и $\omega$ от истинной аномалии начальной орбиты $\nu_1$.

      \item \textbf{Коррекция $r_p$, $r_a$ и $i$.} \\
      \underline{Дано:} \\
      \underline{Начальная орбита:}
      \begin{equation*}
          \begin{array}{r@{\;}c@{\;}l @{\qquad} r@{\;}c@{\;}l}
            r_{p1} & = & 42\;164~\text{км} & \Omega_1 & = & 0^\circ \\
            r_{a1} & = & 42\;164~\text{км} & \omega_1 & = & 0^\circ \\
            i_1    & = & 0.3^\circ         & u_1      & = & 0^\circ \\
          \end{array}
      \end{equation*}
      \underline{Целевые элементы}
      \begin{align*}
            &r_{p2} = 42\;144~\text{км} \\
            &r_{a2} = 42\;184~\text{км} \\
            &i_2 = 0.1^\circ
      \end{align*}
      \underline{Определить:}
       \begin{enumerate}
        \item Импульс скорости $\Delta V$, необходимый для осуществления коррекции.
        \item Эффект на постманевровые параметры $\Omega_2$, $\omega_2$ и $\nu_2$.
      \end{enumerate}
      \underline{Построить} график зависимости хар. скорости манёвра коррекции $r_p$, $r_a$ и $i$ от истинной аномалии начальной орбиты $\nu_1$.
\end{enumerate}

\end{document}
